\scanpage{10}\mspage{\deleted{\unsure{6|8}}7}
% ll. 1–10
\uline{Dém} \deleted{On est ramené par \unsure{principe|procédé} standard} \inserted{La conclusion \inserted{sur $U$} signifie aussi qu'}il existe
un sous-préschéma $U_0$ \deleted{\unsure{$Z$}} de $U$, de prés. finie sur $Y_0$,
tel que $U_0\to U$ soit \uline{surjectif}, \deleted{et tel que} tel que
$U_0\to\deleted{\unsure{X}}Y$ \deleted{soit tel}, soit \deleted{plat} normal \deleted{à \unsure{géom.} \unsure{et} donc loc. de prés finie}
(ce qui implique donc que $U_0$ est réduit, donc $U_0=U_{\text{red}}$), et un ouvert
$V$ de $U_0$, schématiquement dense rel $/Y$, tel que
$V\longrightarrow X$ soit une imm ouverte. Sous cette
forme, on voit qu'on est ramené au cas
où $Y$ est le spectre d'un anneau local ess.
de type fini sur $\mathbb Z$.
\deleted{A prouver qu'il existe une \illegible{1}}

% ll. 11–16
Utilisant la technique \inserted{de constructibilité et} de passage à un \unsure{sch. ét.|échelon|limite},
et le lemme de Hironaka, on voit que l'hypothèse sur la fibre $X_y$ en $x$ est vraie encore aux
pts voisins \inserted{avec la \unsure{m\textsuperscript{me}} dim locale $d$.} Utilisant \inserted{EGA \uline{IV}} 16.1. \deleted{\unsure{et sa}} et la \deleted{\unsure{3\textsuperscript{ème}}}
condition \textcircled{c}, on trouve que $f$ est univ. ouverte \inserted{et équidim} dans
un voisinage de $x$. \deleted{D'autre part, les \unsure{conditions}}
\inserted{On peut donc supposer $f$ univ. ouvert, équidim, à fibres géom. gén. réduites,}
et à \deleted{fibres géom. \unsure{chacunes}} telles que $(\overline X_y)_{\text{réd}}$ soit normal. \deleted{\unsure{Ainsi}} On
\deleted{\unsure{15.3}, on sait que} peut supposer $X$ réduit : prouvons que
$X\to Y$ est normal.
% slanted
\margin{utiliser \textcircled{b} \textcircled{c} \textcircled{d} \textcircled{e}}
% slanted
\margin{\uline{Lemme à \unsure{dégager|dégg.er}.}}
% slanted
\margin{\deleted{On peut \unsure{procéder} par récurrence sur $\dim Y$, et supposer que l'on \unsure{gagne} en dim plus petite}}

% ll. 17–21
Soit $V$ l'ouvert des pts $x$ de $X$ tels que $X_{y=f(x)}$ soit
\uline{réduit} (donc séparable) en $x$ sur $k(y)$. En vertu de \textcircled{a} et
\inserted{EGA \uline{IV}} 15.2.3. c'est aussi le plus grand ouvert sur lequel
$f$ induit un morphisme \deleted{réduit} séparable, et en
vertu de l'hypothèse faite (déduite de \textcircled{d,e}) $V$ est
dense dans $X$ fibre par fibre. \deleted{Si donc $X$ est normal aux \unsure{pts fermés} \unsure{$X_{\text{réd}}$ a les} parties \unsure{réduites},}

% ll. 22–24. Intersecting rulings and text running in opposite directions.
\deleted{Si $\overline X$ est l'adhérence schématique \deleted{\unsure{il sera}} de $V$. Conclusion (que $X$ \inserted{et $Y$} est réduit (qui \unsure{va} se perdre par ext. de la base), \deleted{mais}}
% margin note written sideways (inverted)
\margin{\deleted{\uline{Téléphoner à Coates}}}
% margin note written sideways (inverted)
\margin{\deleted{\unsure{Jean Artin} \unsure{pour 4 heures|pour 1 heure}, voir \unsure{Feitings|fibrés}, \unsure{batterie} \unsure{un} \unsure{cas} \unsure{lost} \unsure{letter}}}

\margin{1.7 ci-dessus}
Le critère valuatif \unsure{nous} ramène, lui, au cas où $Y$ est le sp. d'un anneau de val. discrète,
(par récurr\unsure{ence} !) c'est le cas de Hironaka, OK.
% archive sheet number: 9
